
% ----------------------------------------------------------------------
%                   LATEX TEMPLATE FOR PhD THESIS
% ----------------------------------------------------------------------

% based on Harish Bhanderi's PhD/MPhil template, then Uni Cambridge
% http://www-h.eng.cam.ac.uk/help/tpl/textprocessing/ThesisStyle/
% corrected and extended in 2007 by Jakob Suckale, then MPI-CBG PhD programme
% and made available through OpenWetWare.org - the free biology wiki


%: Style file for Latex
% Most style definitions are in the external file PhDthesisPSnPDF.
% In this template package, it can be found in ./Latex/Classes/
\documentclass[twoside,11pt]{Latex/Classes/PhDthesisPSnPDF}


%: Macro file for Latex
% Macros help you summarise frequently repeated Latex commands.
% Here, they are placed in an external file /Latex/Macros/MacroFile1.tex
% An macro that you may use frequently is the figuremacro (see introduction.tex)
% This file contains macros that can be called up from connected TeX files
% It helps to summarise repeated code, e.g. figure insertion (see below).

% insert a centered figure with caption and description
% parameters 1:filename, 2:title, 3:description and label
\newcommand{\figuremacro}[3]{
	\begin{figure}[htbp]
		\centering
		\includegraphics[width=1\textwidth]{#1}
		\caption[#2]{\textbf{#2} - #3}
		\label{#1}
	\end{figure}
}

% insert a centered figure with caption and description AND WIDTH
% parameters 1:filename, 2:title, 3:description and label, 4: textwidth
% textwidth 1 means as text, 0.5 means half the width of the text
\newcommand{\figuremacroW}[4]{
	\begin{figure}[htbp]
		\centering
		\includegraphics[width=#4\textwidth]{#1}
		\caption[#2]{\textbf{#2} - #3}
		\label{#1}
	\end{figure}
}

% inserts a figure with wrapped around text; only suitable for NARROW figs
% o is for outside on a double paged document; others: l, r, i(inside)
% text and figure will each be half of the document width
% note: long captions often crash with adjacent content; take care
% in general: above 2 macro produce more reliable layout
\newcommand{\figuremacroN}[3]{
	\begin{wrapfigure}{o}{0.5\textwidth}
		\centering
		\includegraphics[width=0.48\textwidth]{#1}
		\caption[#2]{{\small\textbf{#2} - #3}}
		\label{#1}
	\end{wrapfigure}
}

% predefined commands by Harish
\newcommand{\PdfPsText}[2]{
  \ifpdf
     #1
  \else
     #2
  \fi
}

\newcommand{\IncludeGraphicsH}[3]{
  \PdfPsText{\includegraphics[height=#2]{#1}}{\includegraphics[bb = #3, height=#2]{#1}}
}

\newcommand{\IncludeGraphicsW}[3]{
  \PdfPsText{\includegraphics[width=#2]{#1}}{\includegraphics[bb = #3, width=#2]{#1}}
}

\newcommand{\InsertFig}[3]{
  \begin{figure}[!htbp]
    \begin{center}
      \leavevmode
      #1
      \caption{#2}
      \label{#3}
    \end{center}
  \end{figure}
}


%%% Local Variables: 
%%% mode: latex
%%% TeX-master: "~/Documents/LaTeX/CUEDThesisPSnPDF/thesis"
%%% End: 




%: ----------------------------------------------------------------------
%:                  TITLE PAGE: name, degree,..
% ----------------------------------------------------------------------
% below is to generate the title page with crest and author name

%if output to PDF then put the following in PDF header
\ifpdf  
    \pdfinfo { /Title  (PhD and MPhil Thesis Classes)
               /Creator (TeX)
               /Producer (pdfTeX)
               /Author (YourName your@email.net)
               /CreationDate (D:YYYYMMDDhhmmss)  %format D:YYYYMMDDhhmmss
               /ModDate (D:YYYYMMDDhhmm)
               /Subject (xyz)
               /Keywords (add, your, keywords, here) }
    \pdfcatalog { /PageMode (/UseOutlines)
                  /OpenAction (fitbh)  }
\fi


\title{Divergence of an introduced population of the
  Swimbladder-nematode \textit{Anguillicola crassus} - a
  transcriptomic perspective}



% ----------------------------------------------------------------------
% The section below defines www links/email for author and institutions
% They will appear on the title page of the PDF and can be clicked
\ifpdf
  \author{\href{mailto:emanuelheitlinger@gmail.com}{Emanuel Heitlinger}}
  

\cityofbirth{Schw\"abisch Gm\"und} % uncomment this if your university requires this
%  % If city of birth is required, also uncomment 2 sections in PhDthesisPSnPDF
%  % Just search for the "city" and you'll find them.
\collegeordept{\href{http://www.chem-bio.kit.edu/}{Fakult\"at für Chemie und Biowissenschaften}}
\university{\href{http://www.kit.edu/index.php}{Karlsruher Institut
    f\"ur Technologie (KIT) - Universit\"atsbereich}}

  % The crest is a graphics file of the logo of your research institution.
  % Place it in ./0_frontmatter/figures and specify the width
\crest{\includegraphics[width=4cm]{logo}}
  
% If you are not creating a PDF then use the following. The default is PDF.
\else
  \author{Emanuel Heitlinger}
  \cityofbirth{born in Schw\"abisch Gm\"und}
  \collegeordept{Fakult\"at f\"ur Chemie und Biowissenschaften}
  \university{Karlsruher Institut f\"ur Technologie (KIT) - Universit\"atsbereich}
  \crest{\includegraphics[width=4cm]{logo}}
\fi

%\renewcommand{\submittedtext}{change the default text here if needed}
\degree{Dr. rer. nat.}
\degreedate{2012 Januar}


% ----------------------------------------------------------------------
       
% turn of those nasty overfull and underfull hboxes
\hbadness=10000
\hfuzz=50pt


%: --------------------------------------------------------------
%:                  FRONT MATTER: dedications, abstract,..
% --------------------------------------------------------------

\begin{document}

%\language{english}

% sets line spacing
\renewcommand\baselinestretch{1.2}
\baselineskip=18pt plus1pt


%: ----------------------- generate cover page ------------------------

\maketitle  % command to print the title page with above variables


%: ----------------------- cover page back side ------------------------
% Your research institution may require reviewer names, etc.
% This cover back side is required by Dresden Med Fac; uncomment if needed.
%% Is empty the reviewers are on frontpage
\newpage

\newpage


%: ----------------------- abstract ------------------------

% Your institution may have specific regulations if you need an abstract and where it is to be placed in the document. The default here is just after title.


% Thesis Abstract -----------------------------------------------------


%\begin{abstractslong}    %uncommenting this line, gives a different abstract heading
\begin{abstractslong}        %this creates the heading for the abstract page

  The ability to expand into new environments and niches, despite
  being highly adapted to a habitual environment, is a fascinating
  feat of organisms. 30 years ago \textit{Anguillicola crassus} was
  introduced from Asia, where it parasitises \textit{Angilla
    japonica}, to Europe and spread here in the new host species
  \textit{Anguilla anguilla}. Whether and how much phenotypic
  plasticity or rapid adaptation to differential selection are
  contributing to its success in invading new host-populations is a
  question of substantial evolutionary interest.

  Gene regulatory networks, as an important link between genotype
  and phenotype, are thought to play a central role both in the
  response to stress (e.g. from as yet unexperienced environmental
  stressors) and in local adaptation.

  In the present project, differential gene-expression in
  \textit{A. crassus} populations was assessed using next generation
  sequencing on the 454 and Illumina platforms and genetic components
  of differences were isolated in cross-inoculation experiments with
  both Asian and European host-species and parasite populations.

  Several proteases were shown to be under positive selection on the
  sequence level, highlighting this group of enzymes as possible
  targets of an immune-attack on \textit{A. crassus}. On the
  gene-expression level, the extent of heritable change was large in
  comparison to the effect of modification in different
  host-environments. Mitochondrially encoded subunits of the
  respiratory chain and other genes connected to aerobic respiration
  showed divergent expression patterns in European vs. Asian parasite
  populations; cuticle collagen genes showed ``adapted'' patterns of
  expression in present day sympatric host-parasites pairs.
  
  These results identified gene-expression phenotypes, confirming the
  divergence of European \textit{A. crassus} populations. Such
  phenotypes will be more accessible to popluation-genetic analysis
  invetigating selection than complex life history traits.

\end{abstractslong}
%\end{abstractlongs}


% ---------------------------------------------------------------------- 

%%% Local Variables: ***
%%% mode:latex ***
%%% TeX-master: "../thesis.tex"  ***
%%% tex-main-file: "../thesis.tex" ***
%%% End: ***
     

% The original template provides and abstractseparate environment, if your institution requires them to be separate. I think it's easier to print the abstract from the complete thesis by restricting printing to the relevant page.
% \begin{abstractseparate}
%   
% Thesis Abstract -----------------------------------------------------


%\begin{abstractslong}    %uncommenting this line, gives a different abstract heading
\begin{abstractslong}        %this creates the heading for the abstract page

  The ability to expand into new environments and niches, despite
  being highly adapted to a habitual environment, is a fascinating
  feat of organisms. 30 years ago \textit{Anguillicola crassus} was
  introduced from Asia, where it parasitises \textit{Angilla
    japonica}, to Europe and spread here in the new host species
  \textit{Anguilla anguilla}. Whether and how much phenotypic
  plasticity or rapid adaptation to differential selection are
  contributing to its success in invading new host-populations is a
  question of substantial evolutionary interest.

  Gene regulatory networks, as an important link between genotype
  and phenotype, are thought to play a central role both in the
  response to stress (e.g. from as yet unexperienced environmental
  stressors) and in local adaptation.

  In the present project, differential gene-expression in
  \textit{A. crassus} populations was assessed using next generation
  sequencing on the 454 and Illumina platforms and genetic components
  of differences were isolated in cross-inoculation experiments with
  both Asian and European host-species and parasite populations.

  Several proteases were shown to be under positive selection on the
  sequence level, highlighting this group of enzymes as possible
  targets of an immune-attack on \textit{A. crassus}. On the
  gene-expression level, the extent of heritable change was large in
  comparison to the effect of modification in different
  host-environments. Mitochondrially encoded subunits of the
  respiratory chain and other genes connected to aerobic respiration
  showed divergent expression patterns in European vs. Asian parasite
  populations; cuticle collagen genes showed ``adapted'' patterns of
  expression in present day sympatric host-parasites pairs.
  
  These results identified gene-expression phenotypes, confirming the
  divergence of European \textit{A. crassus} populations. Such
  phenotypes will be more accessible to popluation-genetic analysis
  invetigating selection than complex life history traits.

\end{abstractslong}
%\end{abstractlongs}


% ---------------------------------------------------------------------- 

%%% Local Variables: ***
%%% mode:latex ***
%%% TeX-master: "../thesis.tex"  ***
%%% tex-main-file: "../thesis.tex" ***
%%% End: ***
     
% \end{abstractseparate}


%: ----------------------- tie in front matter ------------------------

\frontmatter
% Thesis Dedictation ---------------------------------------------------

\begin{dedication} %this creates the heading for the dedication page

  Natur! [...] Es ist ein ewiges Leben, Werden und Bewegen in ihr, und
  doch rückt sie nicht weiter. Sie verwandelt sich ewig, und ist kein
  Moment Stillestehen in ihr. Fürs Bleiben hat sie keinen Begriff, und
  ihrem Fluch hat sie ans Stillestehen gehängt. Sie ist fest. Ihr
  Tritt ist gemessen, ihre Ausnahmen selten, ihre Gesetze unwandelbar.


  NATURE! [...] Incessant life, development, and movement are in her,
  but she advances not. She changes for ever and ever, and rests not a
  moment. Quietude is inconceivable to her, and she has laid her curse
  upon rest. She is firm. Her steps are measured, her exceptions rare,
  her laws unchangeable.

  J. W. GOETHE translated to English by T. H. HUXLEY for the first
  issue of nature magazine.

  \vspace{6cm}

  To my grandmother Ruth my brother Roman and my wife Silvia

\end{dedication}

% ----------------------------------------------------------------------
% Thesis Acknowledgements ------------------------------------------------


%\begin{acknowledgementslong} %uncommenting this line, gives a different acknowledgements heading
\begin{acknowledgements}      %this creates the heading for the acknowlegments


  I would like to acknowledge the following persons for their
  tremendous help realizing this project (in alphabetical order):

  Mark Blaxter, Stephen Bridgett, Timothee Cezard, Yun-San Han, Sujai
  Kumar, Dominik Laetsch, Jenna Mann, Anna Montazam, Horst
  Taraschewski, Marian Thomson, Urmi Trivedi, Urszula Weclawski,
  Nicola Wrobel and the thousands of unnamed individuals who have
  coded for free software and open source projects used in my
  research.

\end{acknowledgements}
%\end{acknowledgmentslong}

% ------------------------------------------------------------------------





%: ----------------------- contents ------------------------

\setcounter{secnumdepth}{3} % organisational level that receives a numbers
\setcounter{tocdepth}{3}    % print table of contents for level 3
\tableofcontents            % print the table of contents
% levels are: 0 - chapter, 1 - section, 2 - subsection, 3 - subsection


%: ----------------------- list of figures/tables ------------------------

\listoffigures	% print list of figures

\listoftables  % print list of tables


%: ----------------------- glossary ------------------------

% Tie in external source file for definitions: /0_frontmatter/glossary.tex
% Glossary entries can also be defined in the main text. See glossary.tex
% this file is called up by thesis.tex
% content in this file will be fed into the main document

% Glossary entries are defined with the command \nomenclature{1}{2}
% 1 = Entry name, e.g. abbreviation; 2 = Explanation
% You can place all explanations in this separate file or declare them in the middle of the text. Either way they will be collected in the glossary.

% required to print nomenclature name to page header
\markboth{\MakeUppercase{\nomname}}{\MakeUppercase{\nomname}}


% ----------------------- contents from here ------------------------

%  
\nomenclature{SNP}{Single Nucleotide Polymorphism; variation occurring
in a single nucleotide between two closely related homlogous
sequences. Leading to for example to allelic differences within a
population or even the homologous chromosomes in an individual}


\nomenclature{ORF}{Open Reading Frame; a region in a DNA-sequence
begining with a start-codon and not containing a stop-condon. For
example a region within a processed mRNA transcript being transcribed
into a protein}

\nomenclature{dpi}{Days post infection; In infection experiments, a
point in time given in days after an individual has been infected}

\nomenclature{DNA}{Desoxy Ribonucleic Acid; a chemical molecule
  bearing the heritable genetic information in all life on earth}


 

\begin{multicols}{2} % \begin{multicols}{#columns}[header text][space]
\begin{footnotesize} % scriptsize(7) < footnotesize(8) < small (9) < normal (10)

\printnomenclature[1.5cm] % [] = distance between entry and description
\label{nom} % target name for links to glossary

\end{footnotesize}
\end{multicols}



%: --------------------------------------------------------------
%:                  MAIN DOCUMENT SECTION
% --------------------------------------------------------------

% the main text starts here with the introduction, 1st chapter,...
\mainmatter

\renewcommand{\chaptername}{} % uncomment to print only "1" not "Chapter 1"


%: ----------------------- subdocuments ------------------------

% Parts of the thesis are included below. Rename the files as required.
% But take care that the paths match. You can also change the order of appearance by moving the include commands.


% this file is called up by thesis.tex
% content in this file will be fed into the main document

%: ----------------------- introduction file header -----------------------
\chapter{Introduction}

% the code below specifies where the figures are stored
\ifpdf
    \graphicspath{{1_introduction/figures/PNG/}{1_introduction/figures/PDF/}{1_introduction/figures/}}
\else
    \graphicspath{{1_introduction/figures/EPS/}{1_introduction/figures/}}
\fi

% ----------------------------------------------------------------------
%: ----------------------- introduction content ----------------------- 
% ----------------------------------------------------------------------



%: ----------------------- HELP: latex document organisation
% the commands below help you to subdivide and organise your thesis
%    \chapter{}       = level 1, top level
%    \section{}       = level 2
%    \subsection{}    = level 3
%    \subsubsection{} = level 4
% note that everything after the percentage sign is hidden from output



\section{The study organism: \textit{Anguillicola crassus}} 

\subsection{Ecological significance} 

\textit{Anguillicola crassus} Kuwahara, Niimi and Ithakagi 1974
\cite{kuwahara_Niimi_Itagaki_1974, moravec_anguillicoloides} is a
swimbladder nematode naturally parasitizing the Japanese eel
(\textit{Anguilla japonica}) indigenous to East-Asia. After a single
introduction \cite{wielgoss_population_2008} to Germany in the early
1980s \textit{A. crassus} has colonized almost all populations of the
European eel (\textit{Anguilla anguilla})
\cite{kirk_impact_2003}. Since the 1990s populations of the American
eel (\textit{Anguilla rostrata}) have been colonized as novel hosts
\cite{fries_notes:_1996,barse_exotic_1999, barse_swimbladder_2001} and
finally it has been detected in three indigenous \textit{Anguilla}
species on the island of Reunion near Madagascar
\cite{sasal_parasite_2008}.\



In Asia, as well as in the introduced ranges, copepods and ostracods
serve as intermediate hosts of \textit{A. crassus}
\cite{moravec_first_2005}, in which L2 larvae develop to L3 larvae,
infective to the final host. Once ingested by an eel they migrate
through the intestinal wall and via the body cavity into the
swimbladder wall \cite{haenen_effects_1996}, i.a. using a trypsin-like
proteinase\cite{polzer_identification_1993}. In the swimbladder wall
L3 larvae hatch to L4 larvae. After a final moult from L4 to preadult
the parasites inhabit the lumen of the swimbladder, where they
eventually mate. Eggs containing L2 larvae are released via the ductus
pneumaticus into the eels gut and finally into the
water\cite{de_charleroy_life_1990}.\

Within the novel range and hosts, conspicuously elevated prevalences
and intensities of infection occur (reviewed in
\cite{kirk_impact_2003} and \cite{taraschewski_hosts_2007}). These
differences in abundance of \textit{A. crassus} in East Asia compared
to Europe are commonly attributed to the different host-parasite
relations in the final eel host permitting a differential survival of
the larval and the adult parasites
\cite{knopf_differences_2004}. Recently, data from experimental
infections of European eels with \textit{A. crassus} have been
published \cite{fazio_regulation_2008}. They show that the parasite
undergoes (under experimental conditions) a density-dependent
regulation keeping the number of worms within a certain range.\
  

\subsection{Evolutionary significance}

\subsubsection{Divergence of \textit{A. crassus} populations}

\figuremacro{ula_stages.png}{Differences in developmental speed}{data
  courtesy of Urszula Weclawski
  %% \href{http://someurl}{urlanchor<}.
  %% \figuremacroW{ula_stages.png}{Title}{Caption}{0.8}
  % variation of the above macro with a width setting
}

\subsubsection{Interest in \textit{A. crassus based on its
    phylogenetic position in the phylum nematoda}}



\subsection{Functional insights from other nematodes used to formulate
  hypotheses for \textit{A.crassus}}


\section{Advances in sequencing technology enabeling this study}

Recent advances in DNA-sequencing


\figuremacro{sequencing_costs.png}{Falling seqeuncing costs}{Falling into bottomless, Data
  provided by \href{http://www.genome.gov/sequencingcosts/}{National
    Human Genome Research Institute, NHGRI}.}



\subsection{Pyro-sequencing}


\subsection{Illumina-Solexa sequencing}



\begin{table}[htdp]
\centering
\begin{tabular}{ccc} % ccc means 3 columns, all centered; alternatives are l, r

{\bf Gene} & {\bf GeneID} & {\bf Length} \\ 
% & denotes the end of a cell/column, \\ changes to next table row
\hline % draws a line under the column headers

human latexin & 1234 & 14.9 kbps \\
mouse latexin & 2345 & 10.1 kbps \\
rat latexin   & 3456 & 9.6 kbps \\
% Watch out. Every line must have 3 columns = 2x &. 
% Otherwise you will get an error.

\end{tabular}
\caption[title of table]{\textbf{title of table} - Overview of latexin genes.}
% You only need to write the title twice if you don't want it to appear in bold in the list of tables.
\label{latexin_genes} % label for cross-links with \ref{latexin_genes}
\end{table}

% There you go. You already know the most important things.


% ----------------------------------------------------------------------

%%% Local Variables: ***
%%% mode:latex ***
%%% TeX-master: "../thesis.tex"  ***
%%% End: ***
     	% background information
% this file is called up by thesis.tex
% content in this file will be fed into the main document

\chapter{Aims of the project} % top level followed by section, subsection


% ----------------------- contents from here ------------------------

\section{Final aim}

Our ultimate goal is...

\section{Preliminary aims}

There will be several preliminary scientific targets to be accomplished on the way...




% ---------------------------------------------------------------------------
% ----------------------- end of thesis sub-document ------------------------
% ---------------------------------------------------------------------------						% aims of the project

\include{3/XYZ}			
\include{4/XYZ}	
\include{5/XYZ}
\include{6/XYZ}

% this file is called up by thesis.tex
% content in this file will be fed into the main document

\chapter{Discussion} % top level followed by section, subsection


% ----------------------- paths to graphics ------------------------

% change according to folder and file names
\ifpdf
    \graphicspath{{7/figures/PNG/}{7/figures/PDF/}{7/figures/}}
\else
    \graphicspath{{7/figures/EPS/}{7/figures/}}
\fi


% ----------------------- contents from here ------------------------






% ---------------------------------------------------------------------------
% ----------------------- end of thesis sub-document ------------------------
% ---------------------------------------------------------------------------               % discussion of results


% this file is called up by thesis.tex
% content in this file will be fed into the main document

\chapter{Materials \& methods} % top level followed by section, subsection


% ----------------------- paths to graphics ------------------------

% change according to folder and file names
\ifpdf
    \graphicspath{{8/figures/PNG/}{8/figures/PDF/}{8/figures/}}
\else
    \graphicspath{{8/figures/EPS/}{8/figures/}}
\fi

% ----------------------- contents from here ------------------------




 

% ---------------------------------------------------------------------------
%: ----------------------- end of thesis sub-document ------------------------
% ---------------------------------------------------------------------------



 






        % description of lab methods



% --------------------------------------------------------------
%:                  BACK MATTER: appendices, refs,..
% --------------------------------------------------------------

% the back matter: appendix and references close the thesis


%: ----------------------- bibliography ------------------------

% The section below defines how references are listed and formatted
% The default below is 2 columns, small font, complete author names.
% Entries are also linked back to the page number in the text and to
% external URL if provided in the BibTex file.

% PhDbiblio-url2 = names small caps, title bold & hyperlinked, link to page 
\begin{multicols}{2} % \begin{multicols}{ # columns}[ header text][ space]
\begin{tiny} % tiny(5) < scriptsize(7) < footnotesize(8) < small (9)

\bibliographystyle{Latex/Classes/PhDbiblio-url2} % Title is link if provided
\renewcommand{\bibname}{References} % changes the header; default: Bibliography

\bibliography{/home/ele/bibtex/master.bib} % adjust this to fit your BibTex file

\end{tiny}
\end{multicols}

% --------------------------------------------------------------
% Various bibliography styles exit. Replace above style as desired.

% in-text refs: (1) (1; 2)
% ref list: alphabetical; author(s) in small caps; initials last name; page(s)
%\bibliographystyle{Latex/Classes/PhDbiblio-case} % title forced lower case
%\bibliographystyle{Latex/Classes/PhDbiblio-bold} % title as in bibtex but bold
%\bibliographystyle{Latex/Classes/PhDbiblio-url} % bold + www link if provided

%\bibliographystyle{Latex/Classes/jmb} % calls style file jmb.bst
% in-text refs: author (year) without brackets
% ref list: alphabetical; author(s) in normal font; last name, initials; page(s)

%\bibliographystyle{plainnat} % calls style file plainnat.bst
% in-text refs: author (year) without brackets
% (this works with package natbib)


% --------------------------------------------------------------

% according to Dresden med fac summary has to be at the end
%
% Thesis Abstract -----------------------------------------------------


%\begin{abstractslong}    %uncommenting this line, gives a different abstract heading
\begin{abstractslong}        %this creates the heading for the abstract page

  The ability to expand into new environments and niches, despite
  being highly adapted to a habitual environment, is a fascinating
  feat of organisms. 30 years ago \textit{Anguillicola crassus} was
  introduced from Asia, where it parasitises \textit{Angilla
    japonica}, to Europe and spread here in the new host species
  \textit{Anguilla anguilla}. Whether and how much phenotypic
  plasticity or rapid adaptation to differential selection are
  contributing to its success in invading new host-populations is a
  question of substantial evolutionary interest.

  Gene regulatory networks, as an important link between genotype
  and phenotype, are thought to play a central role both in the
  response to stress (e.g. from as yet unexperienced environmental
  stressors) and in local adaptation.

  In the present project, differential gene-expression in
  \textit{A. crassus} populations was assessed using next generation
  sequencing on the 454 and Illumina platforms and genetic components
  of differences were isolated in cross-inoculation experiments with
  both Asian and European host-species and parasite populations.

  Several proteases were shown to be under positive selection on the
  sequence level, highlighting this group of enzymes as possible
  targets of an immune-attack on \textit{A. crassus}. On the
  gene-expression level, the extent of heritable change was large in
  comparison to the effect of modification in different
  host-environments. Mitochondrially encoded subunits of the
  respiratory chain and other genes connected to aerobic respiration
  showed divergent expression patterns in European vs. Asian parasite
  populations; cuticle collagen genes showed ``adapted'' patterns of
  expression in present day sympatric host-parasites pairs.
  
  These results identified gene-expression phenotypes, confirming the
  divergence of European \textit{A. crassus} populations. Such
  phenotypes will be more accessible to popluation-genetic analysis
  invetigating selection than complex life history traits.

\end{abstractslong}
%\end{abstractlongs}


% ---------------------------------------------------------------------- 

%%% Local Variables: ***
%%% mode:latex ***
%%% TeX-master: "../thesis.tex"  ***
%%% tex-main-file: "../thesis.tex" ***
%%% End: ***
     

%: Declaration of originality

% Thesis statement of originality -------------------------------------

% Depending on the regulations of your faculty you may need a declaration like the one below. This specific one is from the medical faculty of the university of Dresden.

\begin{declaration}        %this creates the heading for the declaration page

  I herewith declare that I have produced this paper without the
  prohibited assistance of third parties and without making use of
  aids other than those specified; notions taken over directly or
  indirectly from other sources have been identified as such. This
  paper has not previously been presented in identical or similar form
  to any other German or foreign examination board.

  Chapter \ref{cha:pyro} was in similar form submitted for publication
  to BMC Genomics, in the course of manuscript preparation Mark
  Blaxter edited the text.
  
  Chapter \ref{cha:tra-diff} is in similar intended for publication in
  Plos biology, Mark Blaxter edited parts of the text.

  The thesis work was conducted from May 2008 to December 2011 under
  the supervision of Prof. Dr. Horst Taraschewski at the Karlsruhe
  Institute of Technology and Prof. Mark Blaxter at the University of
  Edinburgh.

\vspace{10mm}

KARLSRUHE, 


\end{declaration}


%%% Local Variables: ***
%%% mode: latex ***
%%% TeX-master: "../thesis.tex" ***
%%% tex-main-file: "../thesis.tex" ***
%%% End: ***

% ----------------------------------------------------------------------



\end{document}
